\documentclass[../../main.tex]{subfiles}% 注意这里的文件路径不能用 ./main.tex ,否则用latexmk编译子文件会报错
\graphicspath{{\subfix{./image/}}} % 指定图片目录,后续可以直接使用图片文件名
% 注意这里的文件路径不能用 ../../image/ ,否则用latexmk编译子文件会报错

% 例如:
% \begin{figure}[H]
% \centering
% \includegraphics[scale=0.3]{图.png}
% \caption{}
% \label{figure:图}
% \end{figure}
% 注意:上述\label{}一定要放在\caption{}之后,否则引用图片序号会只会显示??.

\begin{document}

\section{抽象曲面的曲率}

\begin{definition}
设$(\mathscr{M},g)$为二维抽象曲面，$D$为切向量场的协变微分算子，在容许坐标卡下，定义\textbf{Riemann记号（黎曼曲率张量）}
\begin{align}
R^\delta_{\gamma\beta\alpha}=\frac{\partial\Gamma^\delta_{\gamma\beta}}{\partial u^\alpha}-\frac{\partial\Gamma^\delta_{\gamma\alpha}}{\partial u^\beta}+\Gamma^\mu_{\gamma\beta}\Gamma^\delta_{\mu\alpha}-\Gamma^\mu_{\gamma\alpha}\Gamma^\delta_{\mu\beta}.
\label{eq:8.1}
\end{align}
\end{definition}

\begin{theorem}
Riemann记号$R^\delta_{\gamma\beta\alpha}$在坐标卡变换下按4重线性变换规律变换，其是否为零与坐标卡的选取无关。
\end{theorem}

\begin{proof}
设$(\mathscr{M},g)$的两个容许坐标卡$(U_i,\varphi_i),(U_j,\varphi_j)$满足$U_i\cap U_j\neq\emptyset$，则坐标变换满足
\begin{align}
\frac{\partial}{\partial u_i^\alpha}=\frac{\partial u_j^\xi}{\partial u_i^\alpha}\cdot\frac{\partial}{\partial u_j^\xi},
\label{eq:8.2}
\end{align}
克里斯托费尔符号的变换关系为
\begin{align}
\Gamma^{(i)\delta}_{\gamma\beta}\cdot\frac{\partial u_j^\xi}{\partial u_i^\delta}=\frac{\partial^2 u_j^\xi}{\partial u_i^\gamma\partial u_i^\beta}+\frac{\partial u_j^\zeta}{\partial u_i^\gamma}\cdot\frac{\partial u_j^\eta}{\partial u_i^\beta}\Gamma^{(j)\xi}_{\zeta\eta}.
\label{eq:8.3}
\end{align}
对\eqref{eq:8.3}关于$u_i^\alpha$求导，整理可得
\begin{align}
\left(\frac{\partial\Gamma^{(i)\delta}_{\gamma\beta}}{\partial u_i^\alpha}-\frac{\partial\Gamma^{(i)\delta}_{\gamma\alpha}}{\partial u_i^\beta}\right)\cdot\frac{\partial u_j^\xi}{\partial u_i^\delta}+\Gamma^{(i)\delta}_{\gamma\beta}\cdot\frac{\partial^2 u_j^\xi}{\partial u_i^\delta\partial u_i^\alpha}-\Gamma^{(i)\delta}_{\gamma\alpha}\cdot\frac{\partial^2 u_j^\xi}{\partial u_i^\delta\partial u_i^\beta}=\frac{\partial u_j^\zeta}{\partial u_i^\gamma}\cdot\frac{\partial u_j^\eta}{\partial u_i^\beta}\frac{\partial u_j^\mu}{\partial u_i^\alpha}\left(\frac{\partial\Gamma^{(j)\xi}_{\zeta\eta}}{\partial u_j^\mu}-\frac{\partial\Gamma^{(j)\xi}_{\zeta\eta}}{\partial u_j^\nu}+\Gamma^{(j)\nu}_{\zeta\eta}\Gamma^{(j)\xi}_{\nu\mu}-\Gamma^{(j)\nu}_{\zeta\mu}\Gamma^{(j)\xi}_{\nu\eta}\right).
\label{eq:8.4}
\end{align}
将克里斯托费尔符号变换式代入化简，最终得到
\begin{align}
R^{(i)\delta}_{\gamma\beta\alpha}\cdot\frac{\partial u_j^\xi}{\partial u_i^\delta}=R^{(j)\xi}_{\zeta\eta\mu}\cdot\frac{\partial u_j^\zeta}{\partial u_i^\gamma}\cdot\frac{\partial u_j^\eta}{\partial u_i^\beta}\frac{\partial u_j^\mu}{\partial u_i^\alpha}.
\label{eq:8.5}
\end{align}
因此$R^\delta_{\gamma\beta\alpha}$满足4重线性张量变换规律，其是否为零与坐标选取无关.

\end{proof}

\begin{definition}
定义降指标的黎曼曲率张量与降指标克里斯托费尔符号
\begin{align}
R_{\gamma\delta\alpha\beta}=g_{\delta\eta}R^\eta_{\gamma\alpha\beta},\quad \Gamma_{\delta\alpha\beta}=g_{\delta\eta}\Gamma^\eta_{\alpha\beta},
\label{eq:8.6}
\end{align}
其升指标形式为
$$R^\delta_{\gamma\alpha\beta}=g^{\delta\eta}R_{\gamma\eta\alpha\beta},\quad \Gamma^\delta_{\alpha\beta}=g^{\delta\eta}\Gamma_{\eta\alpha\beta}.$$
进一步定义二维抽象曲面$(\mathscr{M},g)$的\textbf{高斯曲率}
$$K=\frac{R_{1212}}{g_{11}g_{22}-g_{12}g_{21}}.$$
\end{definition}

\begin{remark}
黎曼曲率张量上指标下降时，规定放在下指标的第二个位置，该规定无本质几何意义.
\end{remark}

\begin{proposition}
二维抽象曲面的四降指标黎曼曲率张量满足对称性
\begin{align}
R_{\gamma\delta\beta\alpha}=-R_{\gamma\delta\alpha\beta}=-R_{\delta\gamma\beta\alpha}=R_{\beta\alpha\gamma\delta},
\label{eq:8.8}
\end{align}
且仅含一个实质分量$R_{1212}$.
\end{proposition}

\begin{proof}
由克里斯托费尔符号的度量表达式
$$\Gamma_{\delta\alpha\beta}=\frac12\left(\frac{\partial g_{\alpha\delta}}{\partial u^\beta}+\frac{\partial g_{\delta\beta}}{\partial u^\alpha}-\frac{\partial g_{\alpha\beta}}{\partial u^\delta}\right),$$
代入黎曼曲率张量定义\eqref{eq:8.1}，展开化简得
\begin{align}
R_{\gamma\delta\beta\alpha}&=g_{\delta\mu}\frac{\partial\Gamma^\mu_{\gamma\beta}}{\partial u^\alpha}-g_{\delta\mu}\frac{\partial\Gamma^\mu_{\gamma\alpha}}{\partial u^\beta}+\Gamma^\mu_{\gamma\beta}\Gamma_{\delta\mu\alpha}-\Gamma^\mu_{\gamma\alpha}\Gamma_{\delta\mu\beta}\nonumber\\
&=\frac12\left(\frac{\partial^2 g_{\beta\delta}}{\partial u^\gamma\partial u^\alpha}+\frac{\partial^2 g_{\gamma\alpha}}{\partial u^\delta\partial u^\beta}-\frac{\partial^2 g_{\alpha\delta}}{\partial u^\gamma\partial u^\beta}-\frac{\partial^2 g_{\beta\gamma}}{\partial u^\delta\partial u^\alpha}\right)-\Gamma^\mu_{\gamma\beta}\Gamma_{\mu\delta\beta}+\Gamma^\mu_{\gamma\alpha}\Gamma_{\mu\delta\beta}.
\label{eq:8.7}
\end{align}
由此可直接推导出对称性\eqref{eq:8.8}，二维情形下其余分量均可由$R_{1212}$表示，故仅一个实质分量.

\end{proof}

\begin{theorem}
二维抽象曲面的高斯曲率$K$与容许坐标卡的选取无关，是曲面的内蕴几何量.
\end{theorem}

\begin{proof}
结合黎曼曲率张量的坐标变换关系与度量张量行列式的变换规律
$$g_{11}^{(i)}g_{22}^{(i)}-g_{21}^{(i)}g_{12}^{(i)}=\left(g_{11}^{(j)}g_{22}^{(j)}-g_{21}^{(j)}g_{12}^{(j)}\right)\cdot\left(\frac{\partial u_j^1}{\partial u_i^1}\frac{\partial u_j^2}{\partial u_i^2}-\frac{\partial u_j^2}{\partial u_i^1}\frac{\partial u_j^1}{\partial u_i^2}\right)^2,$$
以及$R_{1212}$的变换关系
\begin{align}
R_{1212}^{(i)}=R_{1212}^{(j)}\cdot\left(\frac{\partial u_j^1}{\partial u_i^1}\frac{\partial u_j^2}{\partial u_i^2}-\frac{\partial u_j^2}{\partial u_i^1}\frac{\partial u_j^1}{\partial u_i^2}\right)^2,
\label{eq:8.9}
\end{align}
两式作商可得
\begin{align}
K=\frac{R_{1212}^{(i)}}{g_{11}^{(i)}g_{22}^{(i)}-g_{21}^{(i)}g_{12}^{(i)}}=\frac{R_{1212}^{(j)}}{g_{11}^{(j)}g_{22}^{(j)}-g_{21}^{(j)}g_{12}^{(j)}}.
\label{eq:8.10}
\end{align}
故$K$与坐标卡选取无关.

\end{proof}

\begin{example}
设$U=\mathbb{R}^2$，度量$g=\frac{4\left((\mathrm{d}u)^2+(\mathrm{d}v)^2\right)}{\left(1+u^2+v^2\right)^2}$，求抽象曲面$S:(U,g)$的高斯曲率.
\end{example}
\begin{solution}
记$u^1=u,\;u^2=v$，则
$$g_{11}=g_{22}=\frac{4}{\left(1+u^2+v^2\right)^2},\quad g_{12}=g_{21}=0,$$
$$g^{11}=g^{22}=\frac{\left(1+u^2+v^2\right)^2}{4},\quad g^{12}=g^{21}=0.$$
计算克里斯托费尔符号：
\begin{align*}
\Gamma^1_{11}&=\frac12\cdot g^{11}\cdot\frac{\partial g_{11}}{\partial u}=-\frac{2u}{1+u^2+v^2},\\
\Gamma^1_{22}&=-\frac12\cdot g^{11}\cdot\frac{\partial g_{22}}{\partial u}=\frac{2u}{1+u^2+v^2},\\
\Gamma^1_{12}&=\Gamma^1_{21}=\frac12\cdot g^{11}\cdot\frac{\partial g_{11}}{\partial v}=-\frac{2v}{1+u^2+v^2},\\
\Gamma^2_{11}&=-\frac12\cdot g^{22}\cdot\frac{\partial g_{11}}{\partial v}=\frac{2v}{1+u^2+v^2},\\
\Gamma^2_{22}&=\frac12\cdot g^{22}\cdot\frac{\partial g_{22}}{\partial v}=-\frac{2v}{1+u^2+v^2},\\
\Gamma^2_{12}&=\Gamma^2_{21}=\frac12\cdot g^{22}\cdot\frac{\partial g_{22}}{\partial u}=-\frac{2u}{1+u^2+v^2}.
\end{align*}
代入黎曼张量公式\eqref{eq:8.7}，得
\begin{align*}
R_{1212}&=g_{22}R^2_{112}\\
&=g_{22}\left(\frac{\partial\Gamma^2_{11}}{\partial v}-\frac{\partial\Gamma^2_{12}}{\partial u}+\Gamma^1_{11}\Gamma^2_{12}+\Gamma^2_{11}\Gamma^2_{22}-\Gamma^1_{12}\Gamma^2_{11}-\Gamma^2_{12}\Gamma^2_{21}\right)\\
&=\frac{4}{\left(1+u^2+v^2\right)^2}\cdot\left[\frac{2}{1+u^2+v^2}-\frac{4v^2}{\left(1+u^2+v^2\right)^2}+\frac{2}{1+u^2+v^2}-\frac{4u^2}{\left(1+u^2+v^2\right)^2}\right]\\
&=\frac{16}{\left(1+u^2+v^2\right)^4}.
\end{align*}
因此
$$K=\frac{R_{1212}}{g_{11}g_{22}-g_{12}g_{21}}=1.$$
该曲面为常曲率$1$的空间，可通过球极投影对应三维欧氏空间中的单位球面.
\end{solution}


\begin{example}
设$U=\{(u,v)\in\mathbb{R}^2:u^2+v^2<1\}$，度量$g=\frac{4\left((\mathrm{d}u)^2+(\mathrm{d}v)^2\right)}{\left(1-u^2-v^2\right)^2}$，求抽象曲面$S:(U,g)$的高斯曲率.
\end{example}
\begin{solution}
由题设，
$$g_{11}=g_{22}=\frac{4}{\left(1-u^2-v^2\right)^2},\quad g_{12}=g_{21}=0,$$
$$g^{11}=g^{22}=\frac{\left(1-u^2-v^2\right)^2}{4},\quad g^{12}=g^{21}=0.$$
计算克里斯托费尔符号：
\begin{align*}
\Gamma^1_{11}&=\frac12\cdot g^{11}\cdot\frac{\partial g_{11}}{\partial u}=\frac{2u}{1-u^2-v^2},\\
\Gamma^1_{22}&=-\frac12\cdot g^{11}\cdot\frac{\partial g_{22}}{\partial u}=-\frac{2u}{1-u^2-v^2},\\
\Gamma^1_{12}&=\Gamma^1_{21}=\frac12\cdot g^{11}\cdot\frac{\partial g_{11}}{\partial v}=\frac{2v}{1-u^2-v^2},\\
\Gamma^2_{11}&=-\frac12\cdot g^{22}\cdot\frac{\partial g_{11}}{\partial v}=-\frac{2v}{1-u^2-v^2},\\
\Gamma^2_{22}&=\frac12\cdot g^{22}\cdot\frac{\partial g_{22}}{\partial v}=\frac{2v}{1-u^2-v^2},\\
\Gamma^2_{12}&=\Gamma^2_{21}=\frac12\cdot g^{22}\cdot\frac{\partial g_{22}}{\partial u}=\frac{2u}{1-u^2-v^2}.
\end{align*}
代入黎曼张量公式\eqref{eq:8.7}，得
\begin{align*}
R_{1212}&=g_{22}R^2_{112}\\
&=g_{22}\left(\frac{\partial\Gamma^2_{11}}{\partial v}-\frac{\partial\Gamma^2_{12}}{\partial u}+\Gamma^1_{11}\Gamma^2_{12}+\Gamma^2_{11}\Gamma^2_{22}-\Gamma^1_{12}\Gamma^2_{11}-\Gamma^2_{12}\Gamma^2_{21}\right)\\
&=\frac{4}{\left(1-u^2-v^2\right)^2}\cdot\left[-\frac{2}{1-u^2-v^2}-\frac{4v^2}{\left(1-u^2-v^2\right)^2}-\frac{2}{1-u^2-v^2}-\frac{4u^2}{\left(1-u^2-v^2\right)^2}\right]\\
&=-\frac{16}{\left(1-u^2-v^2\right)^4}.
\end{align*}
因此
$$K=\frac{R_{1212}}{g_{11}g_{22}-g_{12}g_{21}}=-1.$$
该曲面为常曲率$-1$的空间.

\end{solution}


\end{document}